\section{Methodology}
\label{sec:methodology}

How are candidates generated, scored, statistically tested, and audited against confounds?

\lead{Models, datasets, and caching}
We evaluate five released Base DLMs: DiffuGPT-S, DiffuGPT-M, DiffuLLaMA \citep{gong2025diffullama}, Dream 7B \citep{ye2025dream7b} (adapted from Qwen2.5-7B), and LLaDA 8B \citep{nie2025llada} (trained from scratch, no AR base). On GSM8K \citep{cobbe2021gsm8k}, we split 200 development and 500 held-out test questions; each model generates $N=10$ candidates per question using a fixed four-shot prompt, temperature $T=0.7$, and 256 new tokens across $K\in\{16,32,64\}$ steps, plus a pre-registered random subsample at $K=256$ ($n=185$ Dream, $n=192$ LLaDA, stopped by a pre-declared GPU time limit). We also run a 100-question greedy ($T=0$, $K=256$) fidelity check against published Dream and LLaDA scores. On HumanEval \citep{chen2021humaneval} (164 Python tasks), we split 20 dev and 144 test problems (seed 0) and generate $N=10$ zero-shot candidates at $K=64$, $T=0.7$, and 512 new tokens. For HumanEval+ \citep{liu2023evalplus}, we re-execute the exact same cached candidates under the expanded tests; two test problems where the official Plus canonical solution fails or exceeds our 5\,s CPU cap (\texttt{HumanEval/32}, \texttt{HumanEval/139}) are excluded from both standard and Plus denominators ($n=142$).

\lead{Cut-text boundary and baselines}
Diffusion generations often continue past the final answer line (\texttt{\#\#\#\#~N}), injecting unrelated text that distorts likelihood scorers. We therefore cut every GSM8K candidate immediately after its answer line before AR or decision scoring; keeping the earlier uncut scores exposes the magnitude of this confound. On GSM8K we compare: (i) \emph{first} sample, (ii) \emph{uniform pick} (exact expectation), (iii) \emph{majority vote}, (iv) frozen \emph{AR likelihood} (Qwen2.5-7B, LLaMA-2-7B, Qwen2.5-0.5B, GPT-2, TinyLlama; summed and per-token NLL), (v) diffusion \emph{self-score} on a 200-question gate, (vi) an exploratory \emph{prompted AR choice} baseline giving Qwen2.5-7B the same listwise prompt and reading A--J logits, and (vii) an automated selector search (Arbor) on a separate 300-question split. On HumanEval we compare first, uniform, AR Qwen2.5-7B NLL, a \emph{fixed docstring-example filter} (parsing both \texttt{>>>} doctests and \texttt{f(x)==y}/\texttt{Example:} lines; falls back to candidate 0 if none pass or parse), and an execution-cluster baseline (\emph{MBR-exec}, picking the largest output cluster on docstring inputs).

\lead{Decision-model selectors and pre-registration}
We evaluate the 4B Apache-2.0 single-pass decision model \texttt{pplx-decider-flash-v1.3} (30,720-token context) via OpenRouter's \texttt{/api/alpha/decisions} endpoint (served snapshot \texttt{pplx-decider-flash-v1.3-20261112}). Listwise prompts show the problem and cut candidates A--J (mean 1,544 tokens, max 2,973, zero truncation) and ask which candidate reaches the correct final answer (GSM8K) or correctly implements the docstring and will pass the hidden tests (HumanEval); pointwise prompts score each candidate independently as yes/no (mean 233 tokens). No hidden tests or execution results are ever sent. Primary GSM8K hypotheses (H1: listwise $>$ majority; H2: listwise $>$ AR cut-sum) and HumanEval H4 (listwise $>$ example filter) were committed before scored runs. We declare three deviations: (1) using OpenRouter rather than local 9.6\,GB weights due to download timeout; (2) evaluating H4 on Dream alone after LLaDA collapsed to newline runs at $K=64$ (0/10 on \texttt{HumanEval/0}); and (3) adding ``and will it pass the hidden tests'' to the HumanEval pointwise prompt.

\lead{Execution, statistics, and diagnostic controls}
Code runs in an isolated Python subprocess (\texttt{-I}, 5\,s CPU limit, 2\,GB address space, scrubbed environment, explicit pass sentinel). Because Dream and LLaDA share the same 500 GSM8K questions, naive stacked McNemar double-counts questions; we therefore report \emph{question-clustered bootstrap} 95\% confidence intervals (10,000 resamples over unique questions, seed 0) alongside valid \emph{per-model exact two-sided McNemar} tests. To separate trace-informed selection from answer memorisation or position bias, we run five controls: (1) \emph{judge-alone} (question plus distinct GSM8K candidate answers, reasoning stripped; $n=1,000$); (2) \emph{gold injection} (adding the gold answer as an extra option on the 115 GSM8K items where no candidate is correct); (3) \emph{two-option closed-book} (gold vs.\ a random answer from another GSM8K question; $n=40$); (4) \emph{candidate-order shuffle} (200 GSM8K dev; 50 HumanEval test); and (5) \emph{HumanEval recognition} (canonical vs.\ verified single-token failing mutant on 20 dev + 40 test problems, plus a 20-dev control with all docstrings and examples stripped).
